\documentclass[10pt,a4paper]{amsart}
\usepackage[margin=19mm]{geometry}
\usepackage{amsmath,amssymb,mathtools}
\usepackage[scr=rsfs,cal=euler]{mathalfa}
\usepackage{xfrac}
\usepackage[unicode,hidelinks,bookmarks=false]{hyperref}
\newcommand{\ie}{i.e.\ }
\newcommand{\cf}{cf.\ }
\newcommand{\resp}{respectively\ }
\newcommand{\Subsection}{\subsection}
\providecommand{\nobreakdash}{\nobreak\hbox{-}}
\setlength{\emergencystretch}{2em}
\multlinegap0pt
\usepackage{bm}
\usepackage{bbm}
\usepackage{dsfont}
\usepackage{xspace}
\usepackage{cancel}
\usepackage{mathtools}
\usepackage{amsthm}
\makeatletter
\@ifundefined{theorem}{\newtheorem{theorem}{Theorem}}{}
\makeatother
\newcommand\Cl{\mathcal{C}\ell}
\newcommand{\InMb}{\stackrel{\leftarrow}{\mathcal{L}_{M}}}
\newcommand{\InMf}{\stackrel{\rightarrow}{\mathcal{L}_{M}}}
\renewcommand{\epsilon}{\varepsilon}
\renewcommand{\l}{\lambda}
\title[A generalization of the beam problem: Connection to multi-co]{A generalization of the beam problem: Connection to multi-component Camassa-Holm dynamics}
\author{Richard Beals, Jacek Szmigielski}
\date{Source: arXiv:2509.14639v1 (CC BY 4.0)}
\begin{document}
\maketitle

\noindent\emph{Source and licence.} A generalization of the beam problem: Connection to multi-component Camassa-Holm dynamics. Version arXiv:2509.14639v1 (CC BY 4.0), distributed under the Creative Commons Attribution 4.0 International licence, as recorded in the arXiv licence metadata for this exact version and shown on the abstract page https://arxiv.org/abs/2509.14639v1. Full rights notice: licenses/arxiv-2509.14639-CC-BY-4.0.txt.\par\smallskip

\noindent\textit{Editorial scope.} Counted unit: the Theorem whose source label is \texttt{eq:EB} in the article (source0.tex lines 277--307, source label \texttt{eq:EB}), delivered together with the article's own complete original proof of it (source0.tex lines 308--363, one \texttt{proof} environment of 56 lines). No mathematics is written, repaired or re-derived: statement and proof are the article's own text line for line. Required context retained uncounted: eq:constraints (lines 230--234); eq:t-flow (lines 230--234); eq:Clifdef (lines 255--257). Every definition, display and label of those lines is retained unchanged. Omitted from this package and not redistributed: 1--229, 235--254, 258--276, 364--1137. Nothing inside a retained excerpt is omitted or altered: every retained body is delivered exactly as the article's lines are written, so no omission receipt of any kind accompanies this package. Citations: the retained excerpts carry no citation command at all, so no bibliography entry of the article is transcribed and no reference is redistributed. Reference adaptation: none was needed and none was made; every reference inside the delivered unit resolves inside it, so no unresolved reference or citation is produced. Printed slips kept literal: no printed word, symbol, hypothesis or number of the counted statement or of its proof was altered. Layout and class: the article's own class is \texttt{article} with packages including inputenc, fontenc, microtype, fourier, tikz, caption, amsmath, amsthm, amssymb, mathscinet, pdfpages, empheq, float, subfig; the wrapper is the lane's pinned \texttt{amsart} editorial wrapper at 10pt on A4, which restates the theorem-like environments the retained text needs and adds the article's own macro definitions that the retained text needs (\texttt{\textbackslash Cl}, \texttt{\textbackslash InMb}, \texttt{\textbackslash InMf}, \texttt{\textbackslash epsilon}, \texttt{\textbackslash l}), with extra wrapper packages (bm, bbm, dsfont, xspace, cancel, mathtools). The delivered numbering is the unit's own. Assets: no retained span contains a figure environment, a graphics-inclusion command, a PDF-page inclusion, a table or a TikZ picture; no image file is copied into this package, the package declares no asset dependency and no image file is redistributed with it. Licence: version arXiv:2509.14639v1 of this article is distributed under the Creative Commons Attribution 4.0 International licence, as recorded in the arXiv licence metadata for this exact version and shown on the abstract page https://arxiv.org/abs/2509.14639v1. A full rights notice is delivered at licenses/arxiv-2509.14639-CC-BY-4.0.txt. Characters: every retained excerpt is pure ASCII, so the delivered engine, XeLaTeX with its default Latin Modern text font, reports no missing character.
\begin{align}
\lambda M_t&=-\frac{b_{xxx}}{2}+\frac{\lambda}{2}\big(\InMf b+b\InMb+
[a,M]\big), \label{eq:t-flow}\\
0&=a_x+\lambda[b,M]. \label{eq:constraints} 
\end{align}

\begin{equation} \label{eq:Clifdef}
w_1w_2+w_2w_1=2B(w_1,w_2), \qquad w_1, w_2 \in W, 
\end{equation}

\begin{theorem}
Let us consider the Dirichlet Matrix String problem with values in the Clifford algebra $\Cl(W)$
\begin{align} 
&D_x^2 \Phi=\lambda M \Phi, \qquad   -1<x<1,  \\
\Phi(-1, \l)=0,  \qquad &D_x\Phi(-1,\l)=1, \qquad \text{ and \hspace{0.1cm} $\Phi(+1,\lambda)$ is not invertible}, 
\end{align} 
for $M\in \Cl_1(W)$ and the family of isospectral deformations given by equations \eqref{eq:t-flow} and 
\eqref{eq:constraints}.  Then, there exists a solution to equations \eqref{eq:t-flow} and 
\eqref{eq:constraints}, of the form  
\begin{equation} \label{eq:EB}
b(x,t; \lambda)=b_0(x,t)+\frac{b_{-1}(x,t)}{\lambda}, \quad a(x,t; \lambda)=a(x,t), 
\end{equation} 
with $b_0(x,t) \in \Cl_0(W), \, b_{-1}(x,t)\in \Cl_1(W)$ and $a(x,t)\in \Cl_2(W)$.  More explicitly, if 
\begin{align}
&b_{-1}(x,t)=\beta(x) C=\beta(x)\sum_{\mu=1}^d c_\mu e_\mu,\qquad b_0(x,t)=u(x,t),\\
&      a(x,t)=\sum_{\mu<\nu} a_{\mu\nu }(x,t) e_\mu e_\nu, \hspace{3cm}
\end{align}
where $\beta(x)=1-x^2,$ $C$ is a fixed element in $\Cl_1(W)$, and whenever $\mu<\nu$, 
\begin{equation} \label{eq:vx,uxxx}
 \quad a_{\mu\nu,x}(x,t)=2\beta(x)\begin{vmatrix}m_\mu(x,t)&c_\mu\\ m_\nu(x,t)&c_\nu \end{vmatrix}, \quad u_{xxx}(x,t)=2(\InMf \beta(x), C), 
\end{equation} 
then $M(x,t)$ evolves according to 
\begin{equation} \label{eq:Cliffbeam}
M_t(x,t)=\InMf u(x,t)+\frac12[a,M], 
\end{equation} 
or, in components, 
\begin{equation} \label{eq:Cliffbeamcomp}
m_{\mu,t}(x,t)=\mathcal{L}_{m_\mu(x,t)}u(x,t)+\sum_{\nu=1}^d  a_{\mu\nu}(x,t)\epsilon_\nu m_\nu(x,t),   
\end{equation} 
 where  now $a_{\mu\nu}$ is the skew-symmetric extension of the $a_{\mu \nu}$ in \autoref{eq:vx,uxxx}.   
\end{theorem} 

\begin{proof} 
Using equations \eqref{eq:t-flow}, \eqref{eq:constraints} and \eqref{eq:EB} we get
\begin{equation} 
M_t=\frac12\big(\InMf b_0+b_0\InMb+[a,M]\big) \label{eq:Mt}, 
\end{equation}
and the constraints 
\begin{subequations}
\begin{equation}
b_{0,xxx}=\InMf b_{-1}+b_{-1}\InMb, \label{eq:b0}
\end{equation}
\begin{equation}
b_{-1,xxx}=0, \label{eq:b-1}
\end{equation}
\begin{equation}
[b_0, M]=0, \label{eq:b0M}
\end{equation}
\begin{equation}
a_{x}=[M,b_{-1}].  \label{eq:a0x}
\end{equation}
\end{subequations}

Setting $b_0=u(x,t)$ resolves the constraints \eqref{eq:b0M}.  Likewise, if we set $b_{-1}=\beta(x) C=\beta(x)\sum_{\mu=1}^d c_\mu e_\mu$ then \eqref{eq:b-1} holds.  
In the next step, we compute $\InMf b_{-1}+b_{-1}\InMb$ and $[M,b_{-1}]$ to show that \autoref{eq:vx,uxxx} hold.  Indeed, using \eqref{eq:Clifdef}, we obtain
\begin{equation} 
\begin{split} 
&\InMf b_{-1}+b_{-1}\InMb=(\beta(x)(M(x,t)C+CM(x,t)))_x+\beta_{x}(x)(M(x, t)C+CM(x,t))=\\
 &2B((M(x,t)\beta(x))_x,C)+2B(M(x,t)\beta_x(x),C)=2B(\InMf \beta(x),C).  
 \end{split}
\end{equation} 
Likewise, 
\begin{equation} 
\begin{split}
&[M,b_{-1}]=\beta(x)\sum_{\mu,\nu}m_\mu(x,t)c_\nu[e_\mu,e_\nu]=\\
&\beta(x)\sum_{\mu<\nu}(m_\mu(x,t)c_\nu-m_\mu(x,t) c_\mu)[e_\mu,e_\nu]=2\beta(x)\sum_{\mu<\nu} \begin{vmatrix} m_\mu(x,t)&c_\mu\\m_\nu(x,t)&c_\nu \end{vmatrix} e_\mu e_\nu. 
\end{split} 
\end{equation}
These two formulas, together with \eqref{eq:b0} and \eqref{eq:a0x},  prove \eqref{eq:vx,uxxx}.  Finally, we must compute $\frac12 [a,M]$.  
To this end, we note that 
\begin{equation} 
[e_\mu e_\nu, e_\rho]=2\delta_{\nu \rho} \epsilon_\rho e_\mu-2\delta_{\mu \rho}\epsilon_\rho  e_\nu.  
\end{equation} 
Thus 
\begin{equation} 
\begin{split}
&\frac12[a,M]=\frac12 \sum_{\mu<\nu} \sum_{\rho} a_{\mu \nu}(x,t)m_\rho(x,t)[e_\mu e_\nu, e_\rho]=\\
&\frac12 \big( \sum_{\mu<\nu} a_{\mu \nu}(x,t) 2m_\nu(x,t) \epsilon_\nu e_\mu-\frac12 \sum_{\mu<\nu}a_{\mu \nu}(x,t) 2m_\mu (x,t) \epsilon _\mu  e_\nu \big)=\\
&\sum_{\mu< \nu} a_{\mu \nu }(x,t) m_\nu(x,t)  \epsilon_\nu e_\mu- \sum_{\mu<\nu}a_{\mu \nu }(x,t) m_\mu (x,t) \epsilon_\mu  e_\nu=
\sum_\mu \sum_\nu  a_{\mu \nu}(x,t) m_\nu (x,t) \epsilon_\nu  e_\mu, 
\end{split}
\end{equation} 
with the proviso that $a_{\mu \nu}$, which originally is defined only for $\mu<\nu$,  is extended to a skew-symmetric tensor $a_{\mu \nu}(x,t)$ for all $\mu, \nu$.  Finally, the choice of $\beta(x)=1-x^2$ is dictated by two conditions: 
\begin{enumerate} 
\item $b_{-1},xxx=0$; 
\item $b_{-1}(-1)=b_{-1}(1)=0$  is needed to preserve the Dirichlet boundary conditions.  
\end{enumerate} 
\end{proof} 

\end{document}
